\documentclass[11pt]{article}
\usepackage[margin=1in,top=0.9in,bottom=0.9in]{geometry}
\usepackage{amsmath,amssymb,amsthm,mathtools}
\usepackage[T1]{fontenc}
\usepackage{lmodern,microtype,xurl}
\usepackage[colorlinks=true,linkcolor=blue,citecolor=blue,urlcolor=blue]{hyperref}
\newtheorem{theorem}{Theorem}[section]
\newtheorem{lemma}[theorem]{Lemma}
\newtheorem{proposition}[theorem]{Proposition}
\newtheorem{corollary}[theorem]{Corollary}
\theoremstyle{remark}\newtheorem{remark}[theorem]{Remark}
\newcommand{\R}{\mathbb R}
\newcommand{\E}{\mathbb E}
\newcommand{\HS}{\mathrm{HS}}
\newcommand{\KL}{\operatorname{KL}}
\newcommand{\poly}{\operatorname{poly}}
\newcommand{\PNC}{\mathsf{PNC}}
\newcommand{\Lpoly}{\mathsf L/\mathrm{poly}}
\newcommand{\norm}[1]{\left\lVert #1\right\rVert}
\title{Arithmetic formulas for Gaussian-controlled tanh feature paths}
\author{ChatGPT-6 Astra \and Kaarel H\"anni \and Dmitry Vaintrob}
\date{25 October 2026}
\begin{document}
\maketitle
\begin{abstract}
We study globally margined Boolean classifiers represented by Gaussian
variational feature paths with tanh activation. The Gaussian controls may
have nonzero means, arbitrary covariances, varying query diagonals, and
unrestricted real precision. If their relative entropies are $D_1,\ldots,D_H$,
the classifier has rational arithmetic sign formulas, followed by a majority,
with resources polynomial in
\[
 (d+H+\norm v+2)\exp\!\left[C\sum_{j=1}^H
                   (\sqrt{D_j}+D_j)\right].
\]
Consequently logarithmic entropy occupation gives nonuniform $\PNC^1$
circuits, hence logarithmic space with polynomial advice. The construction
uses normalized Gaussian Hermite features, degree-dependent averaging,
rare nonlinear branches, and an artificial variance floor that removes
hidden assumptions about the precision of importance coefficients.
The Gaussian scalar baseline is strictly larger than the unrestricted
one. Near-baseline consequences are therefore stated relative to the
Gaussian class's own baseline.
This is a theorem about Gaussian variational controls. The unrestricted
critical Bayesian posterior permits non-Gaussian row tilts; no replacement
of those tilts by Gaussian controls is proved here.
\end{abstract}

\section{Model and statement}

Let $d\ge1$, $H\ge1$, $X_d=\{-1,1\}^d$, and
$h_0(x)=(1,x)/\sqrt d$.
A feature state is a map from $X_d$ into a real separable Hilbert space.
At stage $j$, choose a Gaussian row law with mean $m_j$ and covariance
$\Sigma_j>0$, relative to the standard isonormal Gaussian on the preceding
feature span. The next kernel is
\begin{equation}\label{eq:kernel}
 Q_j(x,y)=\E\bigl[\tanh(\langle G_j,h_{j-1}(x)\rangle)
                       \tanh(\langle G_j,h_{j-1}(y)\rangle)\bigr],
 \qquad G_j\sim N(m_j,\Sigma_j).
\end{equation}
Any Hilbert feature realization of this kernel may be used in the next
stage. The control cost is
\begin{equation}\label{eq:entropy}
 D_j=\frac12\left[\norm{m_j}^2+
            \operatorname{tr}(\Sigma_j-I-\log\Sigma_j)\right].
\end{equation}
Only the finite span of the $2^d$ query states matters. We extend a control
by mean zero and covariance identity on its orthogonal complement.
This preserves both the query law and its relative entropy. Coordinates
in the resulting fixed countable feature bases may be arbitrary real
numbers; they need not be sparsely supported or bounded away from zero.

The final output is $z(x)=\langle v,h_H(x)\rangle$. Throughout the main
statement, $f(x)z(x)\ge1$ on the cube. A different positive margin is
handled by scaling $v$.

\begin{theorem}\label{thm:main}
There are universal constants $C,c$ with the following property. Every
Gaussian-controlled path above has a classifier that is a majority of
signs of rational arithmetic formulas whose total size and coefficient
bit lengths are at most $R^c$, where
\[
 R=(d+H+\norm v+2)\exp\!\left[C\sum_j(\sqrt{D_j}+D_j)\right].
\]
Their majority agrees with $f$ on every Boolean input. In particular, whenever
$R$ is polynomial in an ambient input resource, $f$ belongs to nonuniform
$\PNC^1$, and therefore to $\Lpoly$.
\end{theorem}

The arithmetic-formula conclusion is stronger and more precise than a
claim about a real arithmetic computation with shared intermediate states.
All expressions below are scalar trees. A repeated occurrence of an
already sampled child is charged its full size; distinct child calls are
independent. Random choices select an expression independently of the
queried input. No algorithm for finding the nonuniform constants or for
sampling them efficiently is asserted.
Identically zero controls or corrections are omitted. In particular,
a zero branch probability means that no correction is constructed and
no division by that probability is performed.

\paragraph{What is represented.}
The path specifies kernels and a Hilbert-space readout; it does not specify
a finite-width neural draw or a posterior probability. Gaussian means and
covariance changes are charged by~\eqref{eq:entropy}. We first compile
these changes and the Gaussian tanh map into random scalar formulas,
then select and round finitely many formulas on the whole cube. The last
section separately studies the scalar baseline appropriate to this
Gaussian-restricted class. Here $\mathsf{PNC}^1$ means signs of
polynomial-size, logarithmic-depth integer arithmetic circuits in the
nonuniform sense; the selected circuits serve as advice.

\section{Scalar-oracle calculus}

An oracle family $X_i(x)$ has mean vector $u(x)$, pointwise trace moment
$\sum_i\E X_i(x)^2\le V$, and common expected scalar-tree size at most $T$.
We charge $T\ge1$. The coordinate index set may be countable. Trace bounds
are pointwise in $x$; we do not interchange a sum with a supremum over $x$.

\begin{lemma}[Linear and quadratic sampling]\label{lem:sampling}
For a Hilbert--Schmidt matrix $K$, each coordinate of $Ku$ has an unbiased
scalar oracle of expected size $O(T)$, with the sum of its second moments
at most $\norm K_\HS^2 V$. For a Hilbert--Schmidt matrix $B$, the scalar
$u^TBu$ has an unbiased oracle of expected size $O(T)$ and second moment
at most $\norm B_\HS^2V^2$.
\end{lemma}
\begin{proof}
In row $i$, put $b_i=\sum_jK_{ij}^2$, sample $j$ with probability
$K_{ij}^2/b_i$, and return $K_{ij}X_j$ divided by that probability.
Its second moment is $b_i\sum_{j:K_{ij}\ne0}\E X_j^2\le b_iV$.
For the quadratic form, sample $(i,j)$ with probability
$B_{ij}^2/\norm B_\HS^2$ and use independent children $X_iX'_j$.
Tonelli's theorem gives the claimed bound, including countable bases.
\end{proof}

For $A=J+K$, where $J$ is a signed coordinate injection, retain the copied
child and add the sampled $K$ child only on a Bernoulli branch of probability
$p=\norm K_\HS$ when this norm is small. Divide the correction by $p$.
Independence gives
\begin{equation}\label{eq:nearcopy}
 V_{\rm new}\le V+3\norm K_\HS V,
 \qquad T_{\rm new}\le T(1+C\norm K_\HS).
\end{equation}
Indeed the cross term is at most $2\norm K_\HS\norm u^2$ and the new
second-moment term is $\norm K_\HS^2V/p$. On the absent branch we return
the copied formula. Fixed scalar weights are merged along copy chains.
For large $\norm K_\HS$, copy $J$ and sample only $K$, then average
polynomially many copies when the true output mean norm is bounded.
Sampling the whole matrix $J+K$ would not be justified in a countable basis.

An average of two independent oracle families has trace moment
\begin{equation}\label{eq:reset}
 \frac12\left(V+\norm{u(x)}^2\right).
\end{equation}
Thus a bounded-mean family whose moment bound increases by $O(\eta_j)$
per small step can be kept at a fixed bound by $O(1+\sum_j\eta_j)$
two-copy resets. Their cost is $\exp[O(\sum_j\eta_j)]$.

\begin{lemma}[Fixed-dimensional smooth scalar factory]\label{lem:factory}
Let $r\in\R^p$, $p\le2$, have scalar coordinate oracles of uniformly
bounded second moment and expected size $T$. For any compactly supported
$F\in H^4(\R^p)$, there is an unbiased real scalar oracle for $F(r)$ with
second moment at most $C\norm F_{H^4}^2$ and expected size at most $CT$.
The same statement holds for a countable family, with weighted sums of
the second moments bounded by the corresponding weighted sums of
$H^4$ norms.
\end{lemma}
\begin{proof}
Use a fixed frequency density proportional to $(1+\norm\omega^2)^{-4}$.
Construct $W=\sum_{l=1}^p\omega_lX_l$ from independent coordinate
children; its second moment is at most $C\norm\omega^2$. Choose
$K\sim\operatorname{Pois}(\lambda)$ with
$\lambda=1+C\norm\omega^2$, and form
\[
 Z=\prod_{a=1}^K(1+iW_a/\lambda).
\]
The Poisson product identities give $\E Z=e^{i\omega\cdot r}$ and
$\E|Z|^2\le e$. Multiply by the Fourier coefficient divided by the
frequency density and take its real part. Plancherel gives the moment
bound. A balanced real/imaginary expansion uses $O((K+1)^2)$ child
occurrences. The fourth frequency moment is finite in dimensions at most
two for the displayed density, so expected real-tree cost is $O(T)$.
Copies repeated in that expansion use the same already sampled child;
they are not independent redraws. Fixed compact-interval extensions are
used whenever $F$ is specified only on a bounded rectangle.
\end{proof}

\section{Normalized noncentral Gaussian features}

For $b\in\R$, put
\[
 Y_b=\frac{Z+b}{\sqrt{1+b^2}},\quad Z\sim N(0,1),
 \qquad F(t,b)=\E\tanh(\sqrt t\,Y_b)^2.
\]
Then $\E Y_b^2=1$. Expand the normalized scalar feature in orthonormal
probabilists' Hermite polynomials:
\begin{equation}\label{eq:hermite}
 \frac{\tanh(\sqrt t\,Y_b)}{\sqrt{F(t,b)}}
       =\sum_{k\ge0}c_k(t,b)H_k(Z).
\end{equation}
At $t=0$ the continuous extension has
\[
 c_0(0,b)=\frac b{\sqrt{1+b^2}},\qquad
 c_1(0,b)=\frac1{\sqrt{1+b^2}},\qquad c_k(0,b)=0\ (k\ge2).
\]
For a unit vector $u$, the feature map
\begin{equation}\label{eq:fock}
 \mathcal F_{t,b}(u)=\bigoplus_{k\ge0}c_k(t,b)u^{\otimes k}
\end{equation}
has norm one. Using full ordered tensor powers makes its coordinate
description explicit. Its inner products give exactly the normalized
Gaussian tanh kernel, including different $(t,b)$ at different queries.

\begin{lemma}[Coefficient and variance bounds]\label{lem:coefficients}
For all $t>0$ and $b\in\R$,
\begin{equation}\label{eq:weightedtail}
 \sum_{k\ge2}k^2c_k(t,b)^2\le8t^2.
\end{equation}
Mixed derivatives in $t,b$ of every fixed order (in particular through
order five), with two further derivatives in $Z$, of the normalized scalar feature in
\eqref{eq:hermite} are bounded in Gaussian $L^2$ for $0\le t\le1$,
uniformly in $b$. On larger bounded parameter rectangles the bounds are
polynomial in their ranges. Moreover
\begin{equation}\label{eq:envelope}
 F(t,b)\le g(t):=\min\{1,t/(1+t/2)\}.
\end{equation}
\end{lemma}
\begin{proof}
The Gaussian second-derivative identity and $k^2\le2k(k-1)$ for $k\ge2$
give
\[
 \sum_{k\ge2}k^2c_k^2
 \le\frac{2t^2}{(1+b^2)^2 F(t,b)}
             \E\tanh''(\sqrt tY_b)^2\le8t^2,
\]
since $\tanh''(z)^2\le4\tanh(z)^2$.
For the derivative assertion, write
\[
 \tanh(\sqrt tY_b)/\sqrt t=Y_b\psi(tY_b^2),
 \qquad\psi(s)=\tanh(\sqrt s)/\sqrt s.
\]
Every fixed derivative of $\psi$ is bounded on $[0,\infty)$; the Taylor
expansion handles zero and direct differentiation handles infinity.
Every fixed $b$ derivative of $Y_b$ is polynomially bounded in $|Z|$,
uniformly in $b$. Also $F(t,b)/t$ is smooth and bounded away from zero for
$0\le t\le1$, uniformly in $b$: $\E Y_b^2=1$ and uniformly bounded fourth
moments give a uniform small-ball lower bound. Differentiation under the
Gaussian integral now proves the stated bounds. For $t\ge1$, the same
small-ball argument bounds $F(t,b)$ below by a positive constant and
direct differentiation gives polynomial bounds. Finally,
$\tanh^2 z\le z^2/(1+z^2/2)$. For example,
$w(z)=z/\sqrt{1+z^2/2}$ satisfies $w(0)=0$ and
$w'(z)\ge1-w(z)^2$ for $z\ge0$, since
$(1+s)^{-1/2}\ge1-s$ for $s\ge0$; scalar comparison with the tanh
equation proves the inequality. Jensen's inequality for the concave function
$s/(1+s/2)$ and $\E Y_b^2=1$ prove \eqref{eq:envelope}.
\end{proof}

In particular, on a fixed bounded rectangle in $(r,\zeta)$, for small
$a$ and bounded $M$, the functions $c_k(ar,M\zeta)$ have weighted
$H^4$ tail norm
\begin{equation}\label{eq:sobolevtail}
 \sum_{k\ge2}k^2\norm{c_k(a\,\cdot,M\,\cdot)}_{H^4}^2\le Ca^2.
\end{equation}
The two low-degree differences from $c_k(0,M\zeta)$ have $H^4$ norm
$O(a)$, and
\begin{equation}\label{eq:rresidual}
 \norm{F(a r,M\zeta)/g(a)-r}_{H^4}\le Ca.
\end{equation}
These follow from Lemma~\ref{lem:coefficients}, scaling each $t$ derivative
by $a$, and the fact that the differences vanish at $a=0$. For large
parameter ranges, replace the right sides by fixed polynomials in those
ranges. A bounded extension operator into compactly supported $H^4$
functions preserves these weighted estimates.

\section{The path compiler}

We first describe a deterministic-operator version. Let each linear
covariance map be $A_j=J_j+K_j$, with
\[
 \kappa_j=\norm{K_j}_\HS,\quad M_j=\norm{m_j},\quad
 \rho_j=\min\{1,\sigma_{\min}(A_j)\},\quad
 S=\sum_j(\kappa_j+M_j),\quad L=\sum_j\log(1/\rho_j).
\]
The compiler has expected scalar-tree cost
\begin{equation}\label{eq:compilerbudget}
 \poly(H+T_0)\exp[C(S+L)]
\end{equation}
for normalized feature directions. We maintain a unit direction $u(x)$
and a scalar $r(x)\in(0,1]$ such that
\[
 h(x)=\sqrt{\bar q\,r(x)}\,u(x),
\]
where $\bar q$ is deterministic. The combined mean squared norm is at
most two. Common trace moments are kept at a fixed bound using
\eqref{eq:reset}. Initially $r=1$ and the normalized cube coordinates
have bounded deterministic formulas.

\subsection{Control and mean coordinates}

At one stage put
\[
 n=\norm{Au},\qquad\beta=\langle m,u\rangle,\qquad
 u_G=Au/n,\qquad b=\beta/n,\qquad
 c=(1+\kappa)^2+M^2.
\]
The total preactivation second moment is $\bar q r(n^2+\beta^2)$.
Set
\begin{equation}\label{eq:controlstate}
 a=c\bar q,\qquad r'=r(n^2+\beta^2)/c.
\end{equation}
The row at query $x$ has the distribution
$\sqrt{a r'(x)}Y_{b(x)}$, with Gaussian direction $u_G(x)$.
Its leading unit feature at variance zero is
\begin{equation}\label{eq:leading}
 \ell(x)=\frac{(\beta,Au)}{\sqrt{n^2+\beta^2}}
        =\frac{(b,u_G)}{\sqrt{1+b^2}}.
\end{equation}

For a small joint control $\eta=\kappa+M$, the squared-norm corrections
are quadratic forms with Hilbert--Schmidt norm $O(\eta)$:
\[
 A^*A-I=J^*K+K^*J+K^*K,\qquad
 A^*A+mm^*-I=(A^*A-I)+mm^*.
\]
Lemma~\ref{lem:sampling} supplies $O(\eta^2)$-second-moment oracles for
these corrections. Apply Lemma~\ref{lem:factory} after division by
$\eta$ to the smooth function
$(1+\eta s)^{-1/2}-1$, on a fixed interval where its denominator is
bounded away from zero. Its $H^4$ norm is $O(\eta)$. Adding the resulting
normalization products on probability-$\eta$ branches constructs $u_G$
and $\ell$ with moment increase and multiplicative common cost $O(\eta)$.
For the new degree-zero component of $\ell$, create the $\beta$ oracle
only on a probability-$\eta$ branch and divide it by $\eta$ there.
Its trace contribution is $O(M^2/\eta)=O(\eta)$ and its expected
coordinate cost is $1+O(\eta T)$. This avoids a fixed per-layer cost
factor for a newly appended scalar coordinate.
The relative scalar in \eqref{eq:controlstate} is updated in the same way,
using the quadratic form for $n^2+\beta^2-1$.

When $M>0$, a transient scalar helper $\zeta=b/M$ has bounded second
moment and cost $CT$: sample $\langle m,u\rangle/M$ and multiply by the
normalization oracle for $1/n$. It is used only on subsequent rare
nonlinear branches; it is not stored as a new common-cost state that would
introduce a fixed multiplicative factor at each layer. For $M=0$, set
$b=0$ without a helper.
The direction $u_G$ and leading feature $\ell$ are alternative local oracle
families. Their trace budgets are not added as if both were one persistent
state; only $\ell$ is retained as the low-degree output base, while
$u_G$ is called on the charged correction and tensor branches.

For a joint control bounded away from zero, copy $J$ and sample $K$,
then use the lower bound $n\ge\rho$. Smooth compact extensions of the
inverse-square-root functions have Sobolev norms polynomial in
$1+\kappa$, $1+M$, and $1/\rho$. Scalar factories, products, and polynomial
averaging produce the same exact means at a cost polynomial in these
quantities. Its logarithm is $O(\kappa+M+\log(1/\rho))$. The auxiliary
$b$ range is bounded by $M/\rho$ and its rescaling is charged once to
this joint stage.

\subsection{Gaussian features without a tensor variance explosion}

The exact new normalized direction has coordinates
$c_k(a r',b)u_G^{\otimes k}$. For $k\ge2$, average
$m_k=\lceil kV\rceil$ independent copies of each direction-coordinate
oracle before multiplying. Since $\norm{u_G}=1$, their second moments
$W_i(x)$ satisfy
\begin{equation}\label{eq:tensoraverage}
 \left(\sum_iW_i(x)\right)^k
 \le\left(1+\frac{V-1}{m_k}\right)^k\le e.
\end{equation}
Estimate the scalar coefficient independently by
Lemma~\ref{lem:factory}. For small $a$, create this coordinate only with
probability $p_k=a/k^2$, dividing the product by $p_k$ on that branch.
Equation~\eqref{eq:sobolevtail} and \eqref{eq:tensoraverage} give total
tail moment $O(a)$; the common expected coordinate cost is
\[
 1+Cp_k(km_kT+T)\le T(1+Ca).
\]
At degrees zero and one, retain the leading feature \eqref{eq:leading}
and add independent probability-$a$ product corrections for
$c_k(ar',b)-c_k(0,b)$. Their second moments are $O(a^2)$ and their
moment and cost increases are $O(a)$.
The summed leading cross term is nonpositive, since
$c_0(0,b)[c_0(t,b)-c_0(0,b)]+c_1(0,b)[c_1(t,b)-c_1(0,b)]
=\langle Y_b,\phi_{t,b}\rangle-1\le0$ for the two unit Gaussian
functions in \eqref{eq:hermite}.
The scalar $r$ is updated using
\eqref{eq:rresidual} on another probability-$a$ branch. Its new mean is
\begin{equation}\label{eq:qupdate}
 \bar q_{\rm new}=g(a),\qquad
 r_{\rm new}=F(ar',b)/g(a)\in(0,1].
\end{equation}
Thus every small free step costs $1+O(a)$ and has an $O(a)$ reset charge.

For large $a$, use $p_k=\min(1,a/k^2)$, the polynomial Sobolev bounds,
and polynomially many averages. Include the two low degrees and the
scalar variance map directly. This costs polynomially in $1+a$ and
in the parameters of a jointly large control, and restores the combined
trace bound. The logarithmic free-step charge is therefore
$O(a)$ near zero and $O(\log(1+a))$ away from zero.

\subsection{Deterministic variance telescoping}

Put $\ell(a)=\log(a/g(a))$. This is comparable to $a$ near zero and
to $\log(1+a)$ away from zero. Equations \eqref{eq:controlstate} and
\eqref{eq:qupdate} give
\[
 \sum_{j=1}^H\ell(a_j)
 =\log(\bar q_0/\bar q_H)+\sum_j\log c_j.
\]
Since $c_j\ge1$ and $g$ is increasing, the deterministic upper variances
cool no faster than iteration of $g$. After the first step they are at
most one, and $g(q)=q/(1+q/2)$ on $(0,1]$. Hence
\[
 \bar q_H\ge c/(H+1),\qquad
 \sum_j\ell(a_j)\le C\log(H+1)+2\sum_j\log(1+\kappa_j+M_j).
\]
The small-step moment resets and large-step averaging now prove
\eqref{eq:compilerbudget}. All transient helpers are called only within
the explicitly charged branches. No vector subcomputation is shared.

\section{Finite precision without assuming control precision}

We give the precision argument because counting real arithmetic operations
alone would not prove Theorem~\ref{thm:main}. Let $\overline V_i$ denote
the cube average of an oracle's second moment. Remove a zero-$\overline V_i$
oracle; it is zero almost surely at every cube input. Otherwise inject an
independent Rademacher sign:
\begin{equation}\label{eq:floor}
 X_i^+=X_i+\sqrt{\delta\overline V_i}\,\varepsilon_i,
 \qquad
 \operatorname{Var}X_i^+(x)\ge
 \gamma\overline V_i^+,\quad\gamma=\delta/(1+\delta).
\end{equation}
Choose inverse-polynomial $\delta$, inject initially and at stored
intermediates, and restore the floor after averaging resets. Means are
unchanged, trace moments increase by only $O(H\delta)$, and scalar cost
has only an additive polynomial overhead before the previous exponential
budget. Conditional versions of the same floor are used inside scalar
factories below.

Normalize each oracle subtree by its own cube RMS
$a_i=\sqrt{\overline V_i}$. These individual scales can be arbitrarily
small. We place only the resulting coefficient ratios into the fixed
arithmetic tree. The injected-noise node has normalized coefficients
$1/\sqrt{1+\delta}$ and $\sqrt{\delta/(1+\delta)}$.

\paragraph{Linear and sampled product nodes.}
For a near-copy node
$Y=aX_0+(B/p)C_JX_J$ with independent branches, the noise floors imply
\[
 \overline V_Y\ge\gamma\left(a^2\overline V_0+
                 p^{-1}\sum_jp_jC_j^2\overline V_j\right).
\]
Its copied normalized coefficient is at most $\gamma^{-1/2}$, and its
selected normalized coefficient has squared expectation at most
$\gamma^{-1}$. For an independent product, the pointwise second-moment
floors multiply. An $m$-term averaged child has a floor at least
$\gamma/m$ relative to its own cube RMS. Hence a sampled degree-$k$
product has normalized coefficient bounded, on an active branch, by
$p^{-1/2}(m/\gamma)^{k/2}$. Equivalently, its selected squared
coefficient multiplied by the product of the child floors has expectation
at most one. This handles quadratic entry sampling and the independently
estimated coefficient times a tensor product as well.

\paragraph{Fourier factory nodes.}
Condition on frequency and Poisson count $\theta=(\omega,K)$ in
Lemma~\ref{lem:factory}. Let $Z$ be its complex product and put
$s_\theta^2=\E_x\E(|Z|^2\mid\theta)$. If $c$ denotes its Fourier
importance coefficient, replace its conditional real output by
\[
 Y_\theta=\Re(cZ)+\sqrt\delta\,|c|s_\theta\varepsilon.
\]
The added noise charges at most $\delta$ times the same Fourier moment
integral, and the mixed output has a pointwise variance floor. Its
conditional RMS $b_\theta$ satisfies
$b_\theta^2\ge\delta|c|^2s_\theta^2$. Normalize each complex factor by
$t_j=\sqrt{1+r_j^2a_j^2}$, where $r_j$ is its fixed multiplier of the
raw child $X_j$. The normalized factor coefficients are at most one in magnitude,
and the child floors give
\[
 s_\theta^2\ge\gamma^K\prod_jt_j^2.
\]
Thus the conditional output divided by $b_\theta$ has root coefficient
at most $\delta^{-1/2}\gamma^{-K/2}$, plus a noise coefficient at most
one. The mixture normalization contributes $b_\theta/b$, whose squared
expectation is exactly one. This removes possible phase cancellation
without a lower bound on a Fourier coefficient or a frequency cutoff.
For vector frequencies, use the factor
$1+i\sum_l\omega_lX_l/\lambda$ and normalize it by
$t^2=1+\sum_l\omega_l^2a_l^2/\lambda^2$. Its at most three coefficients
are bounded by one. Independent coordinate children give a noise variance
at least $\gamma\sum_l\omega_l^2a_l^2/\lambda^2$, so the same factorwise
lower bound and product proof apply even when their means cancel.

\paragraph{Fixing and rounding a good seed.}
Sample the final linear readout with probabilities proportional to its
squared coordinates. It has second moment at most a constant times
$\norm v^2$. Average polynomially many copies per group, and use $O(d+1)$
independent groups. Chebyshev, a Chernoff bound, and a union bound over the
cube give a realization in which, at every input, at least three quarters
of the group averages are within $1/2$ of the target output. Those groups
have the correct signs and absolute margin at least $1/2$; the other
groups need not have any margin. The expected total formula
size and the expected number of primitive calls are polynomial in $R$.

A size Markov bound and the normalized coefficient second-moment bounds
can be imposed jointly with this success event. On a polynomial-size
realization, every active tensor degree, averaging count, and Poisson
product length is polynomial. Active branches of inverse-superpolynomial
probability may be excluded by an expected-call union bound; equivalently
use the weighted product-coefficient estimate above. Thus, with positive
probability, all normalized coefficients have magnitude at most
$\exp(\poly(R))$. Initial normalized cube leaves are bounded signed
coordinate/constant formulas, and the final root RMS lies between the
margin and a polynomial readout bound. No original control coefficient
or internal RMS scale has to be encoded separately.

Fix this realization. A polynomial-size arithmetic tree with coefficient
magnitudes $\exp(\poly(R))$ has values and coefficient derivatives bounded
by the same form on Boolean inputs. Dyadic rounding to polynomially many
bits therefore changes every group output by less than $1/4$, preserving
the signs of the good three quarters at every input and hence the majority.
Clear positive subtree denominators. Arithmetic-formula
balancing yields polynomial-size logarithmic-depth integer formulas.
If complex gates were retained, first balance using their real-affine
one-hole contexts, then expand real and imaginary parts; this also gives
only polynomial growth. Formula positivity is in $\PNC^1$.
The composition theorem for $\PNC^1$ predicates absorbs the final majority
\cite[Theorem~4.2]{datta}, and $\PNC^1\subseteq\Lpoly$ in the nonuniform
setting follows from \cite[Proposition~2.4]{datta}.

\section{Entropy budgets and the Gaussian scalar baseline}

In the actual Gaussian path, take $A_j=\Sigma_j^{1/2}$ in the current
fixed feature basis. The scalar inequality
$(\sqrt\lambda-1)^2\le\lambda-1-\log\lambda$ gives
\[
 \kappa_j\le\sqrt{2D_j},\qquad M_j\le\sqrt{2D_j}.
\]
A one-dimensional entropy marginal gives
\[
 \log_+\norm{A_j^{-1}}_{\rm op}
    \le C(\sqrt{D_j}+D_j).
\]
Consequently $S+L\le C\sum_j(\sqrt{D_j}+D_j)$. The kernel semantics
are exact throughout: positive covariance square roots act on the existing
feature space, and the noncentral Hermite expansion realizes
\eqref{eq:kernel} including all query-query inner products. The final
physical norm is at most one, so dividing the final feature by its positive
query norm preserves the output sign and cannot reduce a unit margin.
This proves Theorem~\ref{thm:main}.

For $q>0$, let $V_G(q)$ be the infimum over finite scalar paths, including
length zero, of
\[
 \sum_{j=1}^m\KL(\eta_j\Vert N(0,q_{j-1}))+\frac1{8q_m},
 \qquad q_0=q,\quad q_j=\E_{\eta_j}\tanh^2 Z,
\]
where every $\eta_j$ is Gaussian, with arbitrary mean and positive
variance. Write $V(q)$ for the unrestricted scalar value and
$V_*:=\lim_{q\downarrow0}V(q)$.

\begin{proposition}[Different scalar baselines]\label{prop:baselines}
The limit $V_*^G:=\lim_{q\downarrow0}V_G(q)$ exists, and there are
universal $C_0<\infty$ and $\delta_0>0$ such that
\[
 V_*+\delta_0\le V_*^G<2,\qquad
 0\le V_*^G-V_G(q)\le C_0q\quad(0<q\le2).
\]
\end{proposition}
\begin{proof}
Scaling all successive controls proves that $V_G$ is nonincreasing;
Gaussianity and relative entropy are preserved. Let
$F(u)=\E\tanh^2(\sqrt uZ)$. The bound $F(u)\ge u/(1+3u)$ shows
that Gaussian variance inflation from $q$ to $(1+6q)q$, for $q<1/10$,
costs at most $9q^2$ and increases the feature variance by at least
$q^2(3-18q)/(1+3q+18q^2)$. Integrating the cost-to-increment ratio,
and shortening the last step to hit $1/10$, gives total cost including
readout at most
\[
 U=199/200+\log(5/2)<2.
\]
Thus the limit exists and is at most $U$. Likewise
$F(u)\ge u(1-2u)$ shows that inflation to $(1+16q)q$ costs at most
$64q^2$ and advances the state by at least $4q^2$ for $q\le1/128$.
Growing from an arbitrarily small state to $p\le1/128$ costs at most
$17p$. Appending a Gaussian suffix proves
$V_*^G-V_G(p)\le17p$; boundedness extends this to $p\le2$.

To prove strict separation, put $b=1/16$, $a=b/10$, and
\[
 I_p(r)=\sup_{\theta\in\R}\{\theta r-
           \log\E_{N(0,p)}e^{\theta\tanh^2 Z}\}.
\]
For $0<r<1$, the unique entropy-minimizing law with
$\E\tanh^2 Z=r$ has density proportional to
$e^{\theta\tanh^2 Z}$ relative to $N(0,p)$. Its entropy $I_p(r)$
is continuous in $(p,r)$. Consider Gaussian laws $N(m,s)$ with
$p\in[a,b]$, entropy at most two relative to $N(0,p)$, and output
$r=\E_{N(m,s)}\tanh^2 Z\ge b$. Their parameter set is compact:
the entropy bounds $m^2/p$ and confines $s/p$ to a compact subinterval
of $(0,\infty)$. Their outputs are bounded away from one.
On this set the continuous entropy saving
\[
 \KL(N(m,s)\Vert N(0,p))-I_p(r)
\]
has a positive minimum $\delta_0$. Indeed equality would make the
Gaussian log-density ratio a bounded function
$\theta\tanh^2 z-\log\E e^{\theta\tanh^2 Z}$. A quadratic
polynomial can be bounded on $\R$ only if it is constant, forcing
the reference Gaussian and hence $r=F(p)<p\le b$, a contradiction.

Any Gaussian scalar path of cost below two from $q<a$ must cross $b$,
since otherwise its readout costs at least two. At its first crossing,
the preceding state $p$ is at least $a$: a post/pre-state ratio of at
least ten forces entropy at least $(10-1-\log10)/2>2$ by Gaussian
second-moment entropy duality. Replacing that single Gaussian law by
the optimizer defining $I_p(r)$ preserves the post-state and saves at
least $\delta_0$. Hence $V_G(q)\ge V(q)+\delta_0$. Paths of cost
below two approximate $V_G(q)$ by the upper bound $U<2$; taking
$q\downarrow0$ proves the claim.
\end{proof}

\begin{lemma}[Occupation relative to the Gaussian baseline]
\label{lem:gaussianoccupation}
For a globally $1/2$-margined Gaussian path with action
$A=\sum_jD_j+\norm v^2/2\le V_*^G+\epsilon$, $\epsilon\ge0$,
\[
 \sum_j\sqrt{D_j}\le C[1+\log(H+2)+\sqrt{H\epsilon}],
 \qquad \sum_jD_j\le U+\epsilon.
\]
Moreover every such path, without the upper-action hypothesis, satisfies
$A\ge V_*^G-C/H$.
\end{lemma}
\begin{proof}
Let $q_i=\max_x\norm{h_i(x)}^2$. Projection of a Gaussian row law
onto any input is Gaussian, and data processing bounds its entropy by
$D_j$. The scalar terminal readout bound is $1/(8q_H(x))$.
Therefore a scalar Gaussian suffix from an input attaining $q_i$ gives
\[
 S_i:=\sum_{j\le i}D_j\le A-V_G(q_i)\le\epsilon+C_0q_i.
\]
Strict scalar dissipation gives
$D_j\ge q_j-q_{j-1}+cq_{j-1}^2$ for a universal $c>0$.
One direct proof uses
$\log\E e^{\tanh^2(\sqrt pZ)}\le p-cp^2$ on $[0,2]$:
the heat equation makes the inequality strict away from zero, and the
expansion is $p-p^2+O(p^3)$ at zero. For maximum variances, decrease
$c$ until $p-cp^2$ is increasing. These facts also handle $q_0\le2$.

On an interval with pre-states and endpoint at least $u$, entropy $E$
and length $\ell$, divide dissipation by the squared pre-state to get
$\ell\le C(E/u^2+1/u)$. Use last exits from dyadic levels starting
at $\max\{\epsilon,(H+1)^{-1},\min_iq_i\}$. Each later band below
a fixed positive level has entropy $O(u)$ and length $O(1/u)$, so
Cauchy--Schwarz bounds its square-root entropy by a constant.
There are $O(\log(H+2))$ bands. The initial prefix has either entropy
$O(\epsilon+1/H)$ and length at most $H$, or entropy $O(b)$ and
length $O(1/b)$ with $b=\min_iq_i$; its contribution is
$O(1+\sqrt{H\epsilon})$. The remaining high-variance tail has bounded
length. This proves the claim for $\epsilon<1/4$; otherwise apply
Cauchy--Schwarz directly to $\sum D_j\le U+\epsilon$.

For the lower-action assertion it suffices to consider $A\le V_*^G$.
Then $S_i\le C_0q_i$ and $\sum D_j<2$. The same last-exit argument,
starting at $b=\min_iq_i$, bounds the total length by $C/b$.
Thus $b\le C/H$, and its Gaussian suffix costs at least
$V_G(b)\ge V_*^G-C_0b$.
\end{proof}

\begin{corollary}[Near the Gaussian class's own baseline]
\label{cor:critical}
Suppose a globally $1/2$-margined Gaussian-controlled path has action
at most $V_*^G+\epsilon$. Its classifier has $\PNC^1$ resources
polynomial in
\[
 (d+H+\epsilon+2)
 \exp\!\left[C\{1+\log(H+2)+\sqrt{H\epsilon}+\epsilon\}\right].
\]
In particular, at $\epsilon=O(\log^2 R/H)$ with
$d,H\le\poly(R)$ and $H\ge c\log R$, these resources are polynomial
in $R$.
\end{corollary}
\begin{proof}
Scale the margin to one and apply Theorem~\ref{thm:main} and
Lemma~\ref{lem:gaussianoccupation}. The action also bounds the
readout norm by $O(\sqrt{1+\epsilon})$.
\end{proof}

\begin{proposition}[Attainability from the prescribed input]
\label{prop:attainability}
Let $J_H^G(f)$ be the Gaussian-restricted action infimum for global
margin $1/2$, and let $Q_1$ be the first free tanh kernel. Set
\[
 a_f=(f^TQ_1^{-1}f)^{-1},\qquad S_f=1+\log(1/a_f).
\]
There is a universal $C$ such that
\[
 V_*^G-C/H\le J_H^G(f)
 \le V_*^G+C\frac{S_f^2+\log^2(2H)}H\qquad(H\ge CS_f).
\]
One has $S_f\le15(d+1)$, while constants and dictators satisfy
$S_f=O(\log(d+1))$. In particular every fixed target has deep Gaussian
rate limit $V_*^G$, and the own-baseline corollary is nonempty in its
polynomial-resource range.
\end{proposition}
\begin{proof}
The lower bound is Lemma~\ref{lem:gaussianoccupation}. For the upper
bound use the target-component construction of \cite{critical}.
We check specifically that its laws remain Gaussian. The first free
kernel is positive definite and dominates $a_fff^T$. Free cooling,
holding, prescribed growth, and purification use only Gaussian
covariance inflations in this target direction. A smaller tracked
component merely chooses a new decomposition $Q=vff^T+R$, $R\succeq0$;
it changes no law. Representing the reference as
$\sqrt v fZ+G_R$, with independent Gaussian terms, shows by data
processing that each scalar Gaussian modification induces a Gaussian
vector control with no greater entropy, even for singular $R$.

The increasing-path reduction and small-state bridge in
Proposition~\ref{prop:baselines} preserve Gaussianity. The usual dyadic
entropy bound therefore gives a Gaussian scalar suffix from $p$ with
cost at most $V_G(p)+p$ and length $O(1/p)$. Its stable lift rescales
this Gaussian scalar law and adds an independent residual Gaussian,
so it too is admissible. Choosing cooling scale $\eta\asymp S_f/H$
and target scale $p\asymp[S_f+\log(2H)]/H$, the construction has
excess action at most
\[
 C[H\eta^2+S_f\eta+p+p\log(2H)+H^{-4}]
 \le C[S_f^2+\log^2(2H)]/H.
\]
Exact Gaussian holding fills the unused layers before growth. The
first-layer spectral bound gives $a_f\ge e^{-14(d+1)}$; its first
Hermite term gives $a_f\ge c/(d+1)$ for constants and dictators.
These are the asserted bounds on $S_f$.
\end{proof}

Proposition~\ref{prop:baselines} and Lemma~\ref{lem:gaussianoccupation}
show why the choice of baseline matters: the Gaussian subclass with
action at most $V_*+\epsilon$ is empty whenever
$\epsilon<\delta_0-C/H$. A near-$V_*$ Gaussian consequence would
therefore be asymptotically vacuous. The unrestricted posterior admits
all finite-entropy row tilts. Neither a Gaussian replacement theorem nor
a formula compiler for those tilts is proved here, so the original
learning question remains open.

\begin{thebibliography}{9}
\bibitem{datta}
S.~Datta, M.~Mahajan, B.~V.~Raghavendra Rao, M.~Thomas, and H.~Vollmer.
\emph{Counting Classes and the Fine Structure between NC1 and L}.
2010. \url{https://www.cse.iitm.ac.in/~bvrr/conf/mfcs10.pdf}.
\bibitem{critical}
\emph{Critical depth and learning}.
Research manuscript, sections ``Scalar control and target-component
constructions'', ``Short near-optimal scalar paths and residual spectral
control'', and ``The initial target component and the sufficient depth''.
\url{https://github.com/kaarelh/AI-works/tree/main/neural-network-generalization/works-in-progress/critical-depth}.
\end{thebibliography}
\end{document}
